\chapter{Une base de ticks sur des formats ouverts}\label{ch:pl:a-tick-store-on-open-formats}

À 16~h~05, les fichiers de capture de la journée se ferment. À 16~h~40, ils doivent être compactés, triés et interrogeables, car la recherche nocturne
commence à 17~h~00, et un chercheur qui demande la cotation en vigueur avant chacune des transactions de cette semaine ne doit pas remarquer
quels jours proviennent des fichiers écrits pendant la séance et lesquels de l'historique compacté. Une semaine dans la vie d'une telle base
est l'objet de ce chapitre~: cinq jours simulés traversent la capture, une \omterm{def:pl:a-tick-store-on-open-formats:split}{base intrajournalière} écrite au fil de la séance, un compactage à la
clôture, une \omterm{def:pl:a-tick-store-on-open-formats:split}{base historique} de fichiers colonnaires partitionnés, et une couche de requête unique au-dessus des deux. Le quatrième jour, les cotations gagnent un
champ, et la base doit continuer à fonctionner. Rien ici n'exige de base de données propriétaire~: la base est construite à partir des formats ouverts du
chapitre~3, d'un moteur de requête embarqué et d'une bibliothèque de dataframes, avec un fichier plat projeté en mémoire pour les données les plus sollicitées, lu de la même
façon depuis Python, C++ et Rust (\texttt{firm.tickstore}).

\section{Capture en direct~: fichiers en ajout seul}

\begin{definition}[Base de ticks, fichier en ajout seul]\label{def:pl:a-tick-store-on-open-formats:store}
Une \emph{base de ticks}\index{base de ticks} (tick store) est le système qui conserve les \omterm{def:pl:capturing-and-storing-tick-data:tick}{données tick} d'une société de trading et répond aux requêtes sur elles, depuis la séance en cours
jusqu'à l'année la plus ancienne conservée. Un \emph{fichier en ajout seul}\index{fichier en ajout seul} (append-only file) est un fichier dont le writer ajoute des octets
uniquement à sa fin et ne modifie jamais ce qu'il a déjà écrit, de sorte qu'un lecteur peut exploiter chaque enregistrement complet déjà présent pendant que le writer
poursuit son écriture.
\end{definition}

Pendant la séance, la base écrit les données au fur et à mesure de leur arrivée~: dans l'ordre temporel, sans tri, sans compression ni
indexation. Les événements du chapitre~2 --- l'heure de réception suivie de l'événement de 48~octets du feed handler --- sont écrits dans un fichier plat de
enregistrements de largeur fixe~; les cotations et les transactions qui en sont dérivées (le meilleur niveau du carnet après chaque modification, issu du constructeur de carnet (book builder) du livre~13,
et chaque exécution avec son prix) sont écrites dans des flux Arrow IPC, le format colonnaire en mémoire du chapitre~3, sous forme d'une suite de
lots d'enregistrements.

\begin{definition}[Format de fichier plat, version de schéma]\label{def:pl:a-tick-store-on-open-formats:flat}
Un \emph{format de fichier plat}\index{format de fichier plat} (flat file format) range des enregistrements d'une même largeur fixe les uns après les autres, après un en-tête qui
décrit leur disposition, de sorte que l'enregistrement $i$ se trouve à un décalage calculable et qu'il n'est besoin d'analyser aucune donnée pour l'atteindre. La
\emph{version de schéma}\index{version de schéma} (schema version) de l'en-tête identifie la disposition~; un lecteur refuse une version qu'il ne connaît pas.
\end{definition}

Le fichier plat de la base (\cref{fig:pl:a-tick-store-on-open-formats:flat}) commence par un nombre magique, la version, la largeur
des enregistrements et la disposition au format JSON, complétés jusqu'à 64~octets afin que chaque enregistrement qui suit soit aligné de la même manière~; les enregistrements de la version~1 font
56~octets, la version~2 ajoute une plateforme, des indicateurs et un remplissage pour atteindre 64. Le \Cref{lst:pl:a-tick-store-on-open-formats:flat} présente l'écrivain et
le lecteur.

\omcode{code/firm/tickstore/firm_tickstore.py}{55}{83}{Le fichier plat~: un en-tête versionné complété à 64~octets, des enregistrements ajoutés
à la fin, et un lecteur qui projette les enregistrements en mémoire au lieu de les lire.}\label{lst:pl:a-tick-store-on-open-formats:flat}

\begin{omfigure}
\begin{tikzpicture}[font=\scriptsize,f/.style={draw,minimum height=0.55cm,inner sep=2pt,align=center}]
\node[f,fill=gray!15,minimum width=1.0cm] (m) at (0.5,0) {OQTS};
\node[f,fill=gray!15,minimum width=1.3cm,anchor=west,font=\tiny] (v) at (m.east) {v, largeur};
\node[f,fill=gray!10,minimum width=3.1cm,anchor=west,font=\tiny] (s) at (v.east) {schéma JSON, remplissage};
\foreach \i/\c in {0/omDef,1/omMeth,2/omDef,3/omMeth,4/omDef} {\node[f,fill=\c!18,minimum width=1.2cm,anchor=west] at (5.4+\i*1.2,0) {enr.~\i};}
\node at (11.8,0) {\dots};
\draw[decorate,decoration={brace,amplitude=4pt}] (0,0.35) -- (5.4,0.35) node[midway,above=4pt]{en-tête~: multiple de 64~octets};
\draw[decorate,decoration={brace,amplitude=4pt}] (5.4,0.35) -- (11.5,0.35) node[midway,above=4pt]{56 ou 64~octets chacun, ajoutés};
\foreach \x in {5.4,9.4} {\draw[dashed,omThm] (\x,-0.35) -- (\x,-0.9);}
\node[omThm,anchor=north] at (7.4,-0.9) {une page de la projection, lue au premier accès};
\draw[omThm,<->] (5.4,-0.75) -- (9.4,-0.75);
\end{tikzpicture}

\omcaption{Le fichier plat de la base~: un en-tête d'une taille multiple de 64~octets (nombre magique, version, largeur des enregistrements, disposition en JSON),
puis des enregistrements de largeur fixe ajoutés au fur et à mesure de leur arrivée. Un lecteur projette le fichier en mémoire~; le système d'exploitation lit une page depuis le disque la première
fois qu'un enregistrement y est touché.}\label{fig:pl:a-tick-store-on-open-formats:flat}
\end{omfigure}

\begin{remark}[Pourquoi deux formes intrajournalières]\label{rem:pl:a-tick-store-on-open-formats:twoforms}
Le fichier plat est le format le plus rapide à écrire et à parcourir, et la disposition déjà utilisée par le chemin de négociation~; le flux IPC porte un
schéma, des colonnes nullables et de type chaîne, et il se lit dans n'importe quel outil parlant Arrow sans conversion. Une base a besoin des deux, car
ses lecteurs intrajournaliers diffèrent~: un outil de rejeu en C++ veut les enregistrements plats, le notebook d'un chercheur veut une table. Chacune est écrite
une seule fois~; rien n'est converti pendant la séance.
\end{remark}

\section{Compactage de fin de journée}

\begin{definition}[Compactage de fin de journée, clé de tri]\label{def:pl:a-tick-store-on-open-formats:compaction}
Le \emph{compactage de fin de journée}\index{compactage de fin de journée} (end-of-day compaction) réécrit les fichiers intrajournaliers d'une journée dans la disposition historique~: partitionnés,
triés, encodés et compressés pour la lecture. La \emph{clé de tri}\index{clé de tri} (sort key) d'un fichier est la liste des colonnes selon lesquelles ses lignes sont
ordonnées~; les \omterm{def:pl:columnar-formats:pushdown}{cartes de zones} (zone maps) du chapitre~3 élaguent bien sur ses premières colonnes.
\end{definition}

Le compactage de \texttt{firm.tickstore} (\cref{lst:pl:a-tick-store-on-open-formats:compact}) trie chaque table par symbole et
par heure, et écrit un fichier Parquet par partition, sous \texttt{<table>/date=YYYY-MM-DD/symbol=S/}~: les noms des répertoires portent
les valeurs de partition, une convention que les moteurs de requête lisent comme des colonnes. C'est la \cref{met:pl:columnar-formats:layout} du
chapitre~3 rendue concrète~: partitionner par date et par symbole, trier par heure à l'intérieur.

\omcode{code/firm/tickstore/firm_tickstore.py}{147}{161}{Compactage de fin de journée~: tri par symbole puis par heure, et écriture d'un fichier Parquet
par date et par symbole, les valeurs de partition figurant dans les noms de répertoires.}\label{lst:pl:a-tick-store-on-open-formats:compact}

La \Cref{fig:pl:a-tick-store-on-open-formats:sizes} compte les octets de la semaine simulée sous chaque forme. La journée la plus chargée, le lundi,
compte 258\,535 événements~; ses fichiers intrajournaliers occupent 17.5~Mo (les événements plats et les deux flux IPC), ses partitions compactées 6.15~Mo,
2.85 fois moins, avec des événements de 20.0 octets chacun dans leurs partitions. Le compactage de cette journée a pris 0.12~s sur l'ordinateur portable~:
soit environ 2.2 millions d'événements par seconde sur un seul cœur, tri, \omterm{def:pl:capturing-and-storing-tick-data:partition}{partitionnement}, encodage et compression inclus.

\begin{omfigure}
\begin{tikzpicture}
\begin{axis}[width=11.5cm,height=5.8cm,ybar,bar width=8pt,ylabel={Mo},xlabel={jour de négociation (mois-jour)},
  xtick={0,1,2,3,4},xticklabels from table={figdata/platforms/04-a-tick-store-on-open-formats/sizes.csv}{day},
  tick label style={font=\small},label style={font=\small},ymin=0,ymax=16,grid=major,grid style={gray!20},
  legend columns=1,legend style={font=\footnotesize,cells={anchor=west},at={(0.5,-0.3)},anchor=north,draw=none,/tikz/every even column/.append style={column sep=8pt}},
  table/col sep=comma]
\addplot[fill=omDef!45,draw=omDef] table[x=k,y=flat]{figdata/platforms/04-a-tick-store-on-open-formats/sizes.csv};
\addplot[fill=omMeth!45,draw=omMeth] table[x=k,y=ipc]{figdata/platforms/04-a-tick-store-on-open-formats/sizes.csv};
\addplot[fill=omThm!45,draw=omThm] table[x=k,y=parquet]{figdata/platforms/04-a-tick-store-on-open-formats/sizes.csv};
\legend{événements plats (intrajournaliers),flux IPC de cotations et de transactions (intrajournaliers),partitions Parquet (historiques)}
\end{axis}
\end{tikzpicture}

\omcaption{Octets de chaque jour de la semaine simulée sous les formes intrajournalières et, pour les quatre jours compactés, dans les
partitions Parquet historiques (événements, cotations et transactions réunis). Le cinquième jour est celui en cours et n'a pas encore de forme historique~; ses
enregistrements plats font 64~octets, la disposition de la version~2. Données~: \texttt{fig\_tickstore.py}.}\label{fig:pl:a-tick-store-on-open-formats:sizes}
\end{omfigure}

\begin{method}[Compactage sûr]\label{met:pl:a-tick-store-on-open-formats:safe}
(1) Écrire les partitions de la journée dans de nouveaux fichiers placés dans un répertoire de préparation, jamais par-dessus les fichiers intrajournaliers. (2) Les vérifier par rapport aux
données intrajournalières~: même nombre de lignes par table et par symbole, même plage de numéros de séquence avec les mêmes trous, les mêmes
sommes de quantités. (3) Les rendre visibles en une seule étape atomique --- un renommage du répertoire de préparation, ou une nouvelle version d'un
manifeste qui liste les fichiers de la journée. (4) Alors seulement, on supprime les fichiers intrajournaliers, ou on les conserve pour la durée de rétention de la
\omterm{def:pl:capturing-and-storing-tick-data:raw}{capture brute} (chapitre~2).
\end{method}

Un lecteur qui liste un répertoire pendant que le compacteur y écrit voit une demi-journée~; un lecteur passant par un manifeste voit soit l'ancienne
journée, soit la nouvelle. L'étape de vérification est peu coûteuse, car les fichiers intrajournaliers sont encore là, et c'est le dernier moment
où une partition perdue peut être reconstruite sans revenir à la \omterm{def:pl:capturing-and-storing-tick-data:raw}{capture brute}~; le chapitre~7 transforme ces contrôles en règles de qualité
des données de la maison.

\begin{dated}{2026-09}{La conception propriétaire que celui-ci reprend}\label{dat:pl:a-tick-store-on-open-formats:kdb}
La séparation entre intrajournalier et historique est la conception de l'architecture kdb+tick de la base propriétaire kdb+, longtemps répandue dans les banques
et les sociétés de trading~: dans sa documentation (consultée en septembre 2026), un tickerplant écrit les enregistrements entrants dans un journal et les
transmet à une base temps réel, qui conserve la journée en cours en mémoire et l'écrit dans la \omterm{def:pl:a-tick-store-on-open-formats:split}{base historique} en fin de journée~; la \omterm{def:pl:a-tick-store-on-open-formats:split}{base historique}
donne accès à tous les jours précédents. Ce livre enseigne la version à format ouvert de la même conception.
\end{dated}

\section{La base intrajournalière et la base historique}

\begin{definition}[Base intrajournalière, base historique]\label{def:pl:a-tick-store-on-open-formats:split}
La \emph{base intrajournalière}\index{base intrajournalière} conserve les données de la séance en cours sous les formes écrites pendant celle-ci, optimisées pour
l'ajout et pour les lectures récentes~; la \emph{base historique}\index{base historique} conserve chaque journée compactée, optimisée pour
la lecture de plages couvrant plusieurs jours et plusieurs symboles.
\end{definition}

La séparation n'est invisible pour un lecteur que si une seule couche de requête couvre les deux. \texttt{Store} ouvre une base DuckDB embarquée et crée une
vue par table~: les partitions Parquet historiques, lues avec leurs noms de répertoire comme colonnes, et les tables IPC de la journée en cours,
enregistrées en mémoire, réunies par une union qui associe les colonnes par leur nom (\cref{lst:pl:a-tick-store-on-open-formats:store}). Une
requête sur \texttt{quotes} couvre alors le lundi au vendredi, quelle que soit la base qui détient chaque journée.

\omcode{code/firm/tickstore/firm_tickstore.py}{165}{187}{La couche de requête~: une vue par table sur les partitions historiques et la
table intrajournalière de la journée en cours, associées par nom de colonne.}\label{lst:pl:a-tick-store-on-open-formats:store}

\begin{omfigure}
\begin{tikzpicture}[font=\scriptsize,>=Stealth,box/.style={draw,rounded corners=2pt,fill=omDef!10,minimum height=0.7cm,inner sep=3pt,align=center}]
\node[box] (feed) at (0,0.1) {lignes du flux\\(capture)};
\node[box,fill=omMeth!15] (flat) at (3.3,0.7) {événements plats\\(ajout seul)};
\node[box,fill=omMeth!15] (ipc) at (3.3,-0.5) {IPC cotations, transactions\\(lots d'enregistrements)};
\draw[dashed,gray] (1.9,-1.15) rectangle (4.7,1.35);
\node[font=\scriptsize\itshape,anchor=south] at (3.3,1.35) {base intrajournalière};
\node[box,fill=omThm!12] (cmp) at (3.3,-2.4) {compactage\\de fin de journée};
\node[box,fill=omProp!15,text width=2.6cm] (hist) at (7.2,-2.4) {Parquet\\date=D/symbol=S\\triés par temps};
\node[font=\scriptsize\itshape,anchor=north] at (7.2,-3.15) {base historique};
\node[box] (q) at (10.4,0.1) {couche de requête\\(DuckDB, Polars)};
\node[box,fill=gray!10] (u) at (10.4,1.5) {recherche, rejeu};
\draw[->] (feed) -- (flat);
\draw[->] (feed) -- (ipc);
\draw[->] (3.3,-1.15) -- (cmp);
\draw[->] (cmp) -- (hist);
\draw[->] (hist.east) -| (q.south);
\draw[->] (4.7,0.1) -- node[above]{aujourd'hui} (q.west);
\draw[->] (q) -- (u);
\end{tikzpicture}

\omcaption{La \omterm{def:pl:a-tick-store-on-open-formats:store}{base de ticks}~: la capture alimente la \omterm{def:pl:a-tick-store-on-open-formats:split}{base intrajournalière} pendant la séance~; le compactage à la clôture transforme chaque journée en
fichiers colonnaires triés et partitionnés~; une couche de requête unique lit ensemble les partitions historiques et la journée en cours.}\label{fig:pl:a-tick-store-on-open-formats:arch}
\end{omfigure}

\begin{example}[Une fenêtre, deux bases]\label{ex:pl:a-tick-store-on-open-formats:window}
Le nombre de changements de cotation de SIM2 entre 09~h~32 et 09~h~35 pour chaque jour de la semaine donne 12\,291, 996, 1\,388, 1\,197 et
4\,473~: quatre jours proviennent des partitions Parquet, le cinquième de la table IPC intrajournalière, en une seule requête SQL qui ne nomme ni l'une ni l'autre.
La requête a pris 18~ms sur l'ordinateur portable.
\end{example}

\section{Jointures as-of à grande échelle}

La jointure as-of du livre~7, chapitre~3 --- chaque ligne d'une table avec la dernière ligne d'une autre connue à son instant --- est
la requête la plus fréquente de la \omterm{def:pl:a-tick-store-on-open-formats:store}{base de ticks}~: chaque transaction avec la cotation qui prévalait au moment où elle a eu lieu, pour les spreads
effectifs, les signes de transaction et les mark-outs. La base la calcule de trois façons, qui doivent concorder ligne à ligne~: \texttt{ASOF JOIN}
de DuckDB, \texttt{join\_asof} de Polars, et \texttt{firm.pit.asof\_join} du livre~7 appliqué à chaque journée et chaque symbole.

\omcode{code/firm/tickstore/firm_tickstore.py}{190}{202}{La jointure as-of en SQL~: chaque transaction avec la dernière cotation de sa journée et
de son symbole ayant un numéro de séquence antérieur (clé \texttt{seq}), ou à l'instant ou avant son horodatage (clé \texttt{ts}).}\label{lst:pl:a-tick-store-on-open-formats:asof}

La première version du chapitre joignait sur l'horodatage, et les trois moteurs n'étaient pas d'accord. La raison est dans les données~: une
exécution et le changement de cotation qu'elle provoque portent le même horodatage de la bourse, de sorte que «~la dernière cotation à l'instant de la
transaction ou avant~» peut être la cotation que la transaction a elle-même produite, et le moteur ne précise pas laquelle de deux cotations de même horodatage
il renvoie. La cotation qui prévalait \emph{avant} la transaction est la dernière d'un numéro de séquence antérieur, et le numéro de séquence du flux
ordonne totalement les événements.

\begin{proposition}[Jointure sur le numéro de séquence]\label{prop:pl:a-tick-store-on-open-formats:seq}
Si chaque ligne de cotation et de transaction porte le numéro de séquence de l'événement qui l'a produite, la jointure as-of stricte sur le
numéro de séquence renvoie, pour chaque transaction, le meilleur niveau du carnet immédiatement avant l'événement qui l'a exécutée, et le résultat
ne dépend pas de l'algorithme de jointure.
\end{proposition}

\begin{proof}
Les numéros de séquence sont distincts d'un événement à l'autre et croissent dans l'ordre où le carnet les a traités. Les lignes de cotation d'un symbole
dont le numéro de séquence est inférieur à celui de la transaction sont les états de son carnet avant l'exécution~; la dernière d'entre elles est l'état
juste avant celle-ci. Des clés distinctes ne laissent aucune égalité à départager pour un algorithme.
\end{proof}

Sur les 17\,281 transactions de la semaine, les trois moteurs concordent sur chaque ligne avec la jointure sur le numéro de séquence. La jointure sur l'horodatage
renvoie une cotation différente pour 223 d'entre elles, soit 1.3~\%, à chaque fois la cotation que l'exécution venait elle-même de produire~: exactement les
transactions dont un chercheur mesurerait un spread effectif nul ou négatif. Sur ces données, la jointure de DuckDB a pris 16~ms, celle de Polars 45~ms, et
la référence Python 0.17~s (\cref{fig:pl:a-tick-store-on-open-formats:times}).

\begin{omfigure}
\begin{tikzpicture}
\begin{axis}[width=9.8cm,height=6.2cm,xbar,bar width=8pt,xmode=log,log origin=infty,xmin=0.0005,xmax=1,
  xlabel={secondes, meilleur de plusieurs essais (échelle logarithmique)},y dir=reverse,
  ytick={0,1,2,3,4,5,6},yticklabels={compactage du jour le plus chargé,{fenêtre~: SIM2 09~h~32--09~h~35~; cinq jours},jointure as-of~; DuckDB,jointure as-of~; Polars,jointure as-of~; référence (Python),lecture d'un jour~; fichier plat (mmap),lecture d'un jour~; Parquet},
  tick label style={font=\small},label style={font=\small},grid=major,grid style={gray!20},
  point meta=explicit,nodes near coords,nodes near coords style={font=\footnotesize,anchor=west,/pgf/number format/fixed,/pgf/number format/precision=4},
  table/col sep=comma]
\addplot[fill=omDef!45,draw=omDef] table[x=seconds,y expr=\coordindex,meta=seconds]{figdata/platforms/04-a-tick-store-on-open-formats/measured_tickstore.csv};
\end{axis}
\end{tikzpicture}

\omcaption{Durée des principales tâches de la base sur la semaine simulée~: le compactage de la journée la plus chargée (258\,535 événements), la requête
sur la fenêtre de cinq jours, la jointure as-of de 17\,281 transactions par trois moteurs, et un parcours d'une colonne d'une journée depuis le fichier plat
et depuis Parquet. Mesurées sur un ordinateur portable (Intel Core Ultra 7 155H) sous WSL2, un seul fil d'exécution, machine par ailleurs inactive. Données~:
\texttt{bench\_tickstore.py}.}\label{fig:pl:a-tick-store-on-open-formats:times}
\end{omfigure}

\section{Fichiers plats et projection en mémoire}

\begin{definition}[Fichier projeté en mémoire]\label{def:pl:a-tick-store-on-open-formats:mmap}
Un \emph{fichier projeté en mémoire}\index{fichier projeté en mémoire} est un fichier mis en correspondance avec l'espace d'adressage virtuel d'un processus, de sorte que ses octets
sont lus comme de la mémoire~: le système d'exploitation charge chaque page depuis le disque (ou depuis son cache) la première fois qu'elle est touchée, et
les processus qui projettent le même fichier partagent une seule copie de ses pages.
\end{definition}

Un fichier à largeur fixe et une projection en mémoire s'accordent parfaitement~: l'enregistrement $i$ se trouve à une adresse connue, rien n'est analysé, et seules les pages
qu'une requête touche sont lues. Les jeux de données projetés en mémoire du livre~12 (chapitre~23) utilisent le même mécanisme pour les données d'entraînement. Le lecteur de la base
en Python est \texttt{numpy.memmap}~; le lecteur C++20 (\cref{lst:pl:a-tick-store-on-open-formats:cpp}) appelle \texttt{mmap}
lui-même, et le jumeau Rust déclare \texttt{mmap} et \texttt{munmap} depuis la bibliothèque C, puisque le Rust de la série n'utilise aucune
crate externe. Les trois calculent le même résumé par instrument --- nombre d'enregistrements, quantité exécutée, dernier numéro de séquence --- sur les mêmes
jeux d'essai des deux versions. Le parcours d'une colonne des événements du lundi a pris 0.8~ms depuis le fichier plat projeté et 2.1~ms depuis
ses partitions Parquet~: le fichier plat se parcourt plus de deux fois plus vite, mais il occupe 2.8 fois plus de place.

\omcode{code/firm/tickstore/cpp/firm_tickstore.hpp}{29}{62}{Le lecteur C++20~: projeter le fichier, vérifier l'en-tête et lire un enregistrement
à un décalage calculé~; la page qui se trouve sous ce décalage est chargée par le système d'exploitation au premier accès.}\label{lst:pl:a-tick-store-on-open-formats:cpp}

\section{Évolution de schéma}

Le jeudi, la plateforme a ajouté un champ à ses cotations, et les cotations de la base le portent à partir de ce jour. Du lundi au mercredi, aucune
colonne de ce type n'existe.

\begin{definition}[Évolution de schéma]\label{def:pl:a-tick-store-on-open-formats:evolution}
L'\emph{évolution de schéma}\index{evolution de schema@évolution de schéma} est la modification de la disposition d'un jeu de données au fil du temps, selon des règles qui permettent à tout
lecteur de lire toutes les versions~: ici, des champs peuvent être ajoutés en fin de structure avec une valeur par défaut, et ne sont jamais supprimés, renommés, retypés ni réordonnés.
\end{definition}

Selon cette règle, un lecteur de la nouvelle version lit un ancien fichier en remplissant les nouveaux champs avec leurs valeurs par défaut, et un lecteur de l'ancienne
version lit un fichier récent en ignorant les champs qu'il ne connaît pas~; le fichier plat fait les deux grâce à son numéro de version
(\texttt{upgrade}), et la couche de requêtes en mettant les colonnes en correspondance par leur nom (une colonne absente se lit comme nulle). La requête sur la fenêtre
de la même semaine compte la plateforme sur aucune des trois premières journées et sur chaque cotation des deux dernières. La vérification du registre
(\texttt{check\_evolution}) refuse les modifications qui casseraient les anciens lecteurs~: un champ retypé de 32 à 64 bits, un champ
supprimé, deux champs permutés, un champ inséré avant des champs existants.

\begin{method}[Faire évoluer le schéma d'une base de ticks]\label{met:pl:a-tick-store-on-open-formats:evolve}
(1) Ajouter le champ en fin de structure, avec une valeur par défaut qui signifie «~non connu~». (2) Augmenter la version dans l'en-tête et enregistrer la nouvelle
disposition à côté de l'ancienne. (3) Écrire chaque lecteur de façon à accepter toutes les versions enregistrées. (4) Ne jamais réécrire l'historique dans la nouvelle
version, sauf si les anciens lecteurs sont retirés~; si une correction est nécessaire, écrire les journées corrigées comme de nouvelles versions (chapitre~7).
\end{method}

\section{Tutoriel~: une semaine dans la base}\label{tut:pl:a-tick-store-on-open-formats:tut}

\begin{tutorial}
\textbf{Objectif.} Faire passer cinq journées simulées par la capture, la \omterm{def:pl:a-tick-store-on-open-formats:split}{base intrajournalière}, le compactage et la couche de requêtes, avec un changement
de schéma le quatrième jour, et vérifier la jointure as-of de trois façons. \textbf{État final~:} la \cref{fig:pl:a-tick-store-on-open-formats:sizes},
la fenêtre de l'\cref{ex:pl:a-tick-store-on-open-formats:window}, et trois jointures identiques de 17\,281 transactions.
\begin{tutsteps}
\item \textbf{Construire la semaine}~: \texttt{pl\_tickstore.build(root)} capture chaque journée (\texttt{capture\_day}~: les événements vers un
fichier plat, les cotations et les transactions vers des flux IPC, un vidage toutes les 20\,000 événements) et compacte les quatre premières
(le \cref{lst:pl:a-tick-store-on-open-formats:compact}).
\item \textbf{Interroger}~: \texttt{window(root)} compte les cotations de SIM2 pendant les mêmes trois minutes de chaque journée.
\item \textbf{Joindre}~: \texttt{asof\_all(root)} exécute DuckDB, Polars et la référence~; puis relancer DuckDB avec l'horodatage comme clé
et compter les transactions dont la cotation change.
\item \textbf{Projeter}~: construire et exécuter \texttt{firm\_tickstore\_test.cpp} et \texttt{cargo test} dans \texttt{code/firm/tickstore/rust}
sur les jeux d'essai~; comparer avec \texttt{flat\_summary}.
\item \textbf{Mesurer}~: \texttt{bench\_tickstore.py} chronomètre le compactage et les requêtes.
\end{tutsteps}
\textbf{Pour aller plus loin.} Ajouter une troisième \omterm{def:pl:a-tick-store-on-open-formats:flat}{version de schéma} avec un code de condition et lire la semaine entière~; faire écrire au compactage
deux ordres de tri (symbole puis temps, et temps puis symbole) et comparer une requête transversale sur chacun.
\end{tutorial}

\section{Construction~: la base de ticks}\label{bld:pl:a-tick-store-on-open-formats:tickstore}

\begin{build}
\bfield{Rôle} La \omterm{def:pl:a-tick-store-on-open-formats:store}{base de ticks} de la société~: chaque chapitre ultérieur de ce livre qui lit des données de marché les lit ici, le moteur de backtest
du chapitre~11 par sa couche d'accès aux données, les règles de qualité du chapitre~7 sur ses partitions.
\bfield{Interface} \texttt{EVENT\_V1}, \texttt{EVENT\_V2}, \texttt{SCHEMAS}~; \texttt{FlatWriter(path, version)},
\texttt{flat\_open(path) -> (version, memmap)}, \texttt{upgrade}~; \texttt{IpcWriter}, \texttt{ipc\_read}~;
\texttt{quotes\_trades(events, symbols)}~; \texttt{compact(tables, root, date, row\_group\_size)}~;
\texttt{Store(root, intraday, date).sql(q)}~; \texttt{asof\_duckdb}, \texttt{asof\_polars}, \texttt{asof\_reference}~; \texttt{check\_evolution}. C++20
\texttt{firm::tickstore::FlatFile}, \texttt{summarise}~; Rust \texttt{FlatFile::open}, \texttt{get}, \texttt{summarise}.
\bfield{Règles} Les fichiers intrajournaliers sont en ajout seul~; l'en-tête indique la version, et un lecteur refuse une version inconnue~; les
fichiers historiques sont écrits une fois par jour et partitionnés~; chaque cotation et chaque transaction porte le numéro de séquence de son événement~; la
jointure as-of est stricte sur le numéro de séquence~; les schémas ne font que s'agrandir en fin de structure.
\bfield{Tests d'acceptation} \texttt{code/firm/tickstore/tests/}, \texttt{cpp/firm\_tickstore\_test.cpp}, \texttt{rust/}~: fichiers plats des
deux versions, écrits, projetés et refusés lorsqu'ils sont étrangers~; le lecteur Python égal au jeu d'essai que vérifient les jumeaux~; cotations et transactions
d'un carnet établi à la main~; un aller-retour IPC~; la couche de requêtes sur deux journées historiques aux schémas différents et sur la journée courante~;
trois jointures identiques et l'égalité de la jointure sur l'horodatage~; les règles d'évolution.
\bfield{Pour aller plus loin} Un second ordre de tri écrit par le compactage~; une fusion triée des lignes intrajournalières et historiques sans union~;
un compactage d'une journée pendant que la capture de la journée suivante tourne.
\end{build}

\begin{omsources}
\item KX, \emph{kdb+tick} documentation (tickerplant, bases de données temps réel et historiques).
\item DuckDB, \emph{AsOf Join} (documentation)~; Polars, \texttt{DataFrame.join\_asof} (référence de l'API).
\item Apache Arrow, \emph{IPC streaming format} (spécification)~; \texttt{mmap(2)}, page de manuel Linux.
\end{omsources}

\section{Exercices}

\begin{exercise}[$\star$]\label{exo:pl:a-tick-store-on-open-formats:1}
Le fichier plat du lundi fait 14\,478\,152 octets pour 258\,535 enregistrements de version~1. Quelle est la taille de son en-tête, et quelle serait celle du fichier
en version~2~?
\end{exercise}

\begin{exercise}[$\star$]\label{exo:pl:a-tick-store-on-open-formats:2}
Pourquoi un lecteur peut-il utiliser un \omterm{def:pl:a-tick-store-on-open-formats:store}{fichier en ajout seul} pendant que l'écrivain y écrit encore, et que doit-il faire d'un dernier enregistrement incomplet~?
\end{exercise}

\begin{exercise}[$\star$]\label{exo:pl:a-tick-store-on-open-formats:3}
Les fichiers intrajournaliers du lundi occupent 17.5~Mo et ses partitions compactées 6.15~Mo. Quel est le rapport, et d'où vient
l'économie~?
\end{exercise}

\begin{exercise}[$\star\star$]\label{exo:pl:a-tick-store-on-open-formats:4}
Un chercheur calcule des fourchettes effectives à l'aide d'une jointure as-of sur l'horodatage. Sur les données de la semaine, combien de transactions reçoivent la mauvaise
cotation, et dans quel sens leurs fourchettes sont-elles biaisées~?
\end{exercise}

\begin{exercise}[$\star\star$]\label{exo:pl:a-tick-store-on-open-formats:5}
À la cadence de compactage mesurée, combien d'événements le magasin peut-il compacter entre 16~h~05 et 16~h~40 sur un cœur, et
comment ce chiffre se compare-t-il aux 311 milliards de messages par jour prévus pour le flux d'options au chapitre~2~?
\end{exercise}

\begin{exercise}[$\star\star$]\label{exo:pl:a-tick-store-on-open-formats:6}
Un nouveau champ est inséré au milieu de l'enregistrement au lieu d'être ajouté en fin. Que voit un ancien lecteur, et que
signale \texttt{check\_evolution}~?
\end{exercise}

\begin{exercise}[$\star\star\star$]\label{exo:pl:a-tick-store-on-open-formats:7}
\emph{Programmation.} Avec \texttt{flat\_summary} et le lecteur C++, vérifiez le fichier plat du jour courant~: combien d'enregistrements contient chaque
instrument, et quelle quantité a été exécutée sur l'instrument~2~?
\end{exercise}

\begin{exercise}[$\star\star\star$]\label{exo:pl:a-tick-store-on-open-formats:8}
\emph{Trouver la faille.} «~Pour économiser de l'espace, on compacte la journée sur place~: le compacteur lit les fichiers intrajournaliers, écrit les partitions Parquet
par-dessus, et une requête pendant le compactage lit ce qui s'y trouve.~»
\end{exercise}

\section{Problème~: La date limite de 16~h~40}

\begin{problem}[{Problème du week-end --- une semaine dans le magasin de ticks}]\label{pb:pl:a-tick-store-on-open-formats:1}
La semaine simulée du chapitre, son magasin et les mesures de \texttt{bench\_tickstore.py}.

\medskip\noindent\textbf{Partie I --- Intrajournalier.}
\begin{enumerate}
\item Combien d'événements, de cotations et de transactions la semaine contient-elle~?
\item Que contient l'en-tête d'un fichier plat, et pourquoi est-il complété jusqu'à 64~octets~?
\item Pourquoi la \omterm{def:pl:a-tick-store-on-open-formats:split}{base intrajournalière} écrit-elle deux formes, et quels lecteurs utilisent lesquelles~?
\item Quelle est la taille d'un enregistrement dans la version~1 et dans la version~2~?
\item Pourquoi rien n'est-il trié ni compressé pendant la séance~?
\end{enumerate}

\medskip\noindent\textbf{Partie II --- Compactage.}
\begin{enumerate}[resume]
\item Quelle \omterm{def:pl:a-tick-store-on-open-formats:compaction}{clé de tri} et quelle \omterm{def:pl:capturing-and-storing-tick-data:partition}{clé de partition} le compactage utilise-t-il, et pourquoi celles-là~?
\item Combien d'octets le lundi occupe-t-il avant et après le compactage~?
\item Combien de temps le compactage du lundi a-t-il pris sur l'ordinateur portable, et à quel débit~?
\item Combien d'événements les 35 minutes entre 16~h~05 et 16~h~40 pourraient-elles absorber à ce débit sur un cœur~?
\item Que changeriez-vous pour compacter un flux d'options complet dans cette fenêtre~?
\end{enumerate}

\medskip\noindent\textbf{Partie III --- Requêtes.}
\begin{enumerate}[resume]
\item Comment une seule instruction SQL lit-elle les deux bases~?
\item Combien de changements de cotation SIM2 connaît-il entre 09~h~32 et 09~h~35 chaque jour~?
\item Pourquoi la jointure as-of doit-elle utiliser le numéro de séquence plutôt que l'horodatage~?
\item Combien de transactions changent de cotation entre les deux jointures, et quels moteurs s'accordent~?
\item Combien de temps les trois jointures ont-elles pris~?
\end{enumerate}

\medskip\noindent\textbf{Partie IV --- Le verdict.}
\begin{enumerate}[resume]
\item Énoncer le \emph{résultat nommé}~: le débit de compactage et la réduction de taille du jour le plus chargé, les latences des requêtes, et le
plus grand jour que la fenêtre absorbe.
\item Comment les anciennes et les nouvelles cotations coexistent-elles après le changement de schéma de jeudi~?
\item Qu'apporte le fichier plat projeté en mémoire par rapport à Parquet pour le jour courant, et que coûte-t-il~?
\item Quelle partie de cette conception vient des \omterm{def:pl:a-tick-store-on-open-formats:store}{bases de ticks} propriétaires, et laquelle des formats ouverts~?
\item En une phrase~: qu'est-ce qui fait paraître les deux moitiés d'un magasin de ticks comme une seule~?
\end{enumerate}
\end{problem}

\section{Questions d'entretien}

\begin{interviewq}[$\star$ \iqroles{développeur, chercheur}]\label{iq:pl:a-tick-store-on-open-formats:1}
Pourquoi les magasins de ticks séparent-ils le jour courant de l'historique~?
\end{interviewq}

\begin{interviewq}[$\star\star$ \iqroles{développeur, chercheur}]\label{iq:pl:a-tick-store-on-open-formats:2}
Vous joignez les transactions aux cotations à l'instant de la transaction, et certaines fourchettes effectives ressortent négatives. Qu'est-ce qui cloche~?
\end{interviewq}

\begin{interviewq}[$\star\star$ \iqroles{développeur}]\label{iq:pl:a-tick-store-on-open-formats:3}
Qu'apporte la projection d'un fichier en mémoire, et quand est-ce une mauvaise idée~?
\end{interviewq}

\begin{interviewq}[$\star\star$ \iqroles{développeur}]\label{iq:pl:a-tick-store-on-open-formats:4}
Une plateforme ajoute un champ à son flux. Comment modifiez-vous le magasin sans casser cinq années de lecteurs~?
\end{interviewq}

\begin{interviewq}[$\star\star$ \iqroles{développeur}]\label{iq:pl:a-tick-store-on-open-formats:5}
Comment partitionneriez-vous et trieriez-vous une année de cotations pour 8\,000 symboles, et pourquoi~?
\end{interviewq}

\begin{interviewq}[$\star\star\star$ \iqroles{développeur}]\label{iq:pl:a-tick-store-on-open-formats:6}
Concevez un magasin de ticks sur formats ouverts pour une société sans base de données propriétaire~: capture, intrajournalier,
compactage, disposition de l'historique, couche de requêtes, et ce que vous mesureriez en premier.
\end{interviewq}
